% ============================================================================
%  HANDOUT — Un crash course de Dynare
%  Del papel al código: un modelo RBC en Dynare
%  Ventura & Echevarría — Universidad del Pacífico, Macroeconomía II
%
%  Autocontenido: compila con pdflatex (un solo motor), sin biblatex/biber.
%  Preamble adaptado de las Lecture Notes de Macroeconomics II.
% ============================================================================
\documentclass[11pt,a4paper]{report}

% --- CODIFICACIÓN E IDIOMA ---
\usepackage[utf8]{inputenc}
\usepackage[T1]{fontenc}
\usepackage[spanish, es-tabla, es-nodecimaldot]{babel}

% --- TIPOGRAFÍA Y MATEMÁTICAS ---
\usepackage[sc]{mathpazo}
\linespread{1.05}
\usepackage{amsmath, amssymb, amsfonts, mathtools, bm}
\usepackage{siunitx}
\usepackage{amsthm}

% --- DISEÑO DE PÁGINA ---
\usepackage{geometry}
\geometry{a4paper, top=2.5cm, bottom=2.5cm, left=3cm, right=2.5cm}
\usepackage{fancyhdr}
\usepackage{xcolor}

% --- COLORES (idénticos al libro) ---
\definecolor{darkblue}{rgb}{0.0, 0.0, 0.55}
\definecolor{defgreen}{RGB}{235, 250, 235}
\definecolor{defframe}{RGB}{70, 150, 70}
\definecolor{thmblue}{RGB}{228, 238, 255}
\definecolor{thmframe}{RGB}{40, 80, 170}
\definecolor{propblue}{RGB}{233, 242, 255}
\definecolor{propframe}{RGB}{60, 100, 190}
\definecolor{lemmacyan}{RGB}{230, 248, 252}
\definecolor{lemmaframe}{RGB}{20, 130, 150}
\definecolor{corolgray}{RGB}{244, 244, 244}
\definecolor{corolframe}{RGB}{120, 120, 120}
\definecolor{exgray}{RGB}{250, 250, 248}
\definecolor{exframe}{RGB}{150, 150, 140}
\definecolor{remarkamber}{RGB}{255, 251, 228}
\definecolor{remarkframe}{RGB}{200, 130, 0}
\definecolor{assbg}{RGB}{255, 245, 232}
\definecolor{assframe}{RGB}{190, 95, 0}
\definecolor{upblue}{RGB}{15, 35, 85}

% --- TABLAS Y FIGURAS ---
\usepackage{booktabs}
\usepackage{graphicx}
\usepackage{float}
\usepackage{tikz}
\usetikzlibrary{positioning, calc, arrows.meta}
\usepackage[labelfont=bf, font=small]{caption}

% --- CÓDIGO (listings; para los capítulos de código, ya preparado) ---
\usepackage{listings}
\lstdefinestyle{dynare}{
  basicstyle=\ttfamily\footnotesize,
  keywordstyle=\color{darkblue}\bfseries,
  commentstyle=\color{exframe}\itshape,
  stringstyle=\color{lemmaframe},
  showstringspaces=false,
  breaklines=true,
  frame=single,
  rulecolor=\color{gray!40},
  numbers=left, numberstyle=\tiny\color{gray!70},
  xleftmargin=1.5em,
  morekeywords={var,varexo,parameters,model,end,steady_state_model,shocks,
    stoch_simul,steady,check,resid,initval,stderr,order,irf}
}
\lstset{style=dynare}

% --- TCOLORBOX (solo breakable: sin skins) ---
\usepackage[breakable]{tcolorbox}
\tcbset{
  defstyle/.style={breakable, colback=defgreen, colframe=defframe,
    colbacktitle=defframe, coltitle=white, fonttitle=\bfseries\small,
    boxrule=0.5pt, arc=2pt, toptitle=1.5mm, bottomtitle=1.5mm,
    left=3pt, right=3pt, top=2pt, bottom=2pt},
  thmstyle/.style={breakable, colback=thmblue, colframe=thmframe,
    colbacktitle=thmframe, coltitle=white, fonttitle=\bfseries\small,
    boxrule=0.5pt, arc=2pt, toptitle=1.5mm, bottomtitle=1.5mm,
    left=3pt, right=3pt, top=2pt, bottom=2pt},
  propstyle/.style={breakable, colback=propblue, colframe=propframe,
    colbacktitle=propframe, coltitle=white, fonttitle=\bfseries\small,
    boxrule=0.5pt, arc=2pt, toptitle=1.5mm, bottomtitle=1.5mm,
    left=3pt, right=3pt, top=2pt, bottom=2pt},
  examplestyle/.style={breakable, colback=exgray, colframe=exframe,
    colbacktitle=exframe, coltitle=white, fonttitle=\bfseries\small,
    boxrule=0.5pt, arc=2pt, toptitle=1.5mm, bottomtitle=1.5mm,
    left=3pt, right=3pt, top=2pt, bottom=2pt},
  remarkstyle/.style={breakable, colback=remarkamber, colframe=remarkframe,
    colbacktitle=remarkframe, coltitle=white, fonttitle=\bfseries\itshape\small,
    leftrule=4pt, rightrule=0.4pt, toprule=0.4pt, bottomrule=0.4pt, arc=0pt,
    left=5pt, right=3pt, top=2pt, bottom=2pt, toptitle=1.5mm, bottomtitle=1.5mm}
}

% --- CONTADORES ---
\newcounter{tcbdef}[section]
\newcounter{tcbex}[section]
\newcounter{tcbrem}[section]
\renewcommand{\thetcbdef}{\thesection.\arabic{tcbdef}}
\renewcommand{\thetcbex}{\thesection.\arabic{tcbex}}
\renewcommand{\thetcbrem}{\thesection.\arabic{tcbrem}}

% --- ENTORNOS ---
\newenvironment{definition}[1][]{%
  \refstepcounter{tcbdef}%
  \begin{tcolorbox}[defstyle, title={Definición~\thetcbdef%
    \ifx\relax#1\relax\else\ (#1)\fi}]}%
{\end{tcolorbox}}
\newenvironment{example}[1][]{%
  \refstepcounter{tcbex}%
  \begin{tcolorbox}[examplestyle, title={Ejemplo~\thetcbex%
    \ifx\relax#1\relax\else\ (#1)\fi}]}%
{\end{tcolorbox}}
\newenvironment{remark}{%
  \refstepcounter{tcbrem}%
  \begin{tcolorbox}[remarkstyle, title={Observación~\thetcbrem}]}%
{\end{tcolorbox}}

% --- TÍTULOS Y LISTAS ---
\usepackage{titlesec}
\usepackage{tocloft}
\usepackage{enumitem}

% --- HIPERVÍNCULOS ---
\usepackage{hyperref}
\hypersetup{
  colorlinks=true, linkcolor=darkblue, citecolor=darkblue, urlcolor=darkblue,
  pdftitle={Un crash course de Dynare — Del papel al código},
  pdfauthor={Ventura \& Echevarría}}

% --- CABECERAS ---
\pagestyle{fancy}
\fancyhf{}
\fancyhead[L]{\small\itshape Un crash course de Dynare}
\fancyhead[R]{\small\thepage}
\renewcommand{\headrulewidth}{0.5pt}
\fancypagestyle{plain}{\fancyhf{}\fancyfoot[C]{\thepage}%
  \renewcommand{\headrulewidth}{0pt}}

% --- FORMATO DE CAPÍTULOS (estilo del libro) ---
\titleformat{\chapter}[display]
  {\bfseries\Large}
  {\filleft\MakeUppercase{\chaptertitlename}\Huge\thechapter}
  {4ex}
  {\titlerule\vspace{2ex}\filleft}
  [\vspace{2ex}\titlerule]

% --- MACROS ---
\newcommand{\code}[1]{\texttt{#1}}
\newcommand{\Real}{\mathbb{R}}
\newcommand{\E}{\mathbb{E}}

% ============================================================================
\begin{document}

% ===================== PORTADA =====================
\begin{titlepage}
\centering
\vspace*{1.2cm}
\includegraphics[width=0.60\textwidth]{assets/UP_logo.jpg}\\[2.2cm]
{\Large\scshape Repensando la Macroeconomía}\\[0.9cm]
{\color{upblue}\rule{\linewidth}{1.2pt}}\\[0.5cm]
{\Huge\bfseries Un crash course de Dynare}\\[0.5cm]
{\large\itshape Del papel al código: un modelo RBC en Dynare}\\[0.5cm]
{\color{upblue}\rule{\linewidth}{1.2pt}}\\[2.2cm]
{\large Alejandro Ventura \quad\&\quad Marcelo Echevarría}\\[0.5cm]
{\normalsize Universidad del Pacífico}\\[0.2cm]
{\normalsize Semestre 2026-II}\\
\vfill
{\small Material para la sesión práctica (hands-on). Acompaña a las diapositivas del curso.}
\end{titlepage}

\tableofcontents
\thispagestyle{plain}
\clearpage

% ===================== CAPÍTULO 1 =====================
\chapter{Introducción a Dynare}

\begin{quote}\itshape
Un modelo macroeconómico que no puede ser llevado a los datos ni simulado para
generar predicciones cuantitativas es, en el mejor de los casos, un ejercicio
intelectual incompleto. Dynare es la herramienta que cierra esa brecha:
convierte las ecuaciones abstractas de un DSGE en trayectorias observables,
funciones de impulso-respuesta y distribuciones posteriores. Aprender a usarlo
es aprender a \textup{\textit{hablar}} el lenguaje de la macroeconomía aplicada
moderna.
\end{quote}

\bigskip

\section{Motivación: por qué bajar de la pizarra}

A lo largo del curso hemos construido, con lápiz y pizarra, el andamiaje de un
modelo de equilibrio general dinámico: las preferencias del hogar, el problema
de optimización intertemporal, las condiciones de primer orden, el vaciado de
los mercados. Todo ese trabajo culmina, invariablemente, en un sistema de
ecuaciones no lineales que describe las trayectorias de equilibrio de la
economía.

El ejercicio que resolveremos en esta sesión ilustra bien hasta dónde llega el
lápiz. Para el modelo RBC con trabajo que estudiaremos, la pizarra \emph{sí}
entrega una solución cerrada, pero solo bajo el supuesto artificial de
depreciación total ($\delta=1$): en ese caso las horas no reaccionan al choque
de productividad y el modelo colapsa a algo que podemos resolver a mano. En
cuanto relajamos ese supuesto y dejamos que el capital se deprecie de a pocos
($\delta<1$) ---es decir, en cuanto el modelo se vuelve \emph{realista}--- la
solución cerrada desaparece y necesitamos la computadora. Esa frontera, entre
lo que la pizarra resuelve y lo que exige una herramienta numérica, es
precisamente el asunto que nos trae a colación.

\textbf{Dynare} es la plataforma de referencia para cruzar esa frontera. Es
software libre, corre sobre MATLAB u Octave, y lo usan bancos centrales,
ministerios de finanzas y departamentos de economía de todo el mundo. Su
filosofía de diseño es separar dos capas del trabajo del investigador. La
primera es la \emph{descripción económica}: uno escribe en un archivo de texto
(\code{.mod}) las ecuaciones del modelo casi tal como aparecerían en el papel.
La segunda es la \emph{computación}: Dynare interpreta ese archivo y aplica por
nosotros los algoritmos numéricos pertinentes ---descomposición de Schur,
filtros de Kalman, métodos de Monte Carlo--- sin que debamos programarlos. Así,
el investigador concentra su esfuerzo en las decisiones económicas del modelo,
mientras Dynare se ocupa de resolverlo y simularlo numéricamente.

\begin{definition}[Dynare]
Dynare es una plataforma de software de código abierto para el manejo de una
amplia clase de modelos económicos dinámicos, en particular modelos DSGE. Sus
capacidades incluyen: (i)~simulación determinista bajo previsión perfecta;
(ii)~simulación estocástica mediante aproximaciones de Taylor de orden~$n$
alrededor del estado estacionario; (iii)~estimación bayesiana vía filtro de
Kalman y MCMC; y (iv)~análisis de política óptima.
\end{definition}

El objetivo de esta sesión es concreto: tomar el modelo RBC del ejercicio,
escribirlo en Dynare de principio a fin ---declaraciones, ecuaciones, estado
estacionario, calibración--- y usarlo para generar funciones de
impulso-respuesta y momentos. Al terminar, cada quien habrá \emph{corrido} su
propio modelo DSGE, no solo leído sobre él.

\section{De la pizarra al computador: una breve historia}

Conviene detenerse un momento en \emph{de dónde viene} esta forma de hacer
macroeconomía, porque Dynare no cae del cielo: es el último eslabón de una
transformación que tomó dos décadas y que cambió lo que significa ``tener un
modelo''.

\begin{remark}
Dynare no es la única forma de llevar un modelo DSGE al computador. Hoy existe
un ecosistema amplio ---Python (con bibliotecas como QuantEcon, de Sargent y
Stachurski), Julia, R--- y no faltan quienes sostienen que, para muchos
problemas modernos, las bibliotecas científicas de Python ya superan a MATLAB
(véase QuantEcon, \url{https://python-programming.quantecon.org}), además de
integrarse mejor con campos como el aprendizaje automático. En esta sesión
optamos por Dynare, sobre MATLAB u Octave, por tres razones prácticas: su
sintaxis declarativa permite escribir el modelo casi como en el papel, es el
estándar \emph{de facto} en bancos centrales y la academia, y los alumnos de la
UP cuentan con licencia de MATLAB. La elección de plataforma, en todo caso, no
cambia la economía del modelo: lo que aprendamos aquí se traslada a cualquier
otro entorno.
\end{remark}

\subsection{La visión: el modelo como laboratorio}

El punto de quiebre intelectual lo articuló Robert E.\ Lucas~Jr.\ en
1980. Su propuesta fue redefinir qué es una teoría económica: no un conjunto de
afirmaciones sobre la economía real, sino un \emph{sistema artificial} que
podemos construir y manipular. En sus palabras:

\begin{quote}\itshape
``One of the functions of theoretical economics is to provide fully articulated,
artificial economic systems that can serve as laboratories in which policies
that would be prohibitively expensive to experiment with in actual economies
can be tested out at much lower cost.''
\end{quote}
\hfill{\small(Lucas, 1980, p.~696)}

\noindent La idea del modelo como \emph{laboratorio} ---un lugar donde ensayar
políticas que sería demasiado caro o peligroso probar en la economía real---
exige, por supuesto, que ese laboratorio pueda efectivamente ponerse a
funcionar. Y aquí Lucas fue sorprendentemente literal sobre lo que ello
significaba:

\begin{quote}\itshape
``Our task, as I see it \ldots is to write a FORTRAN program that will accept
specific economic policy rules as `input' and will generate as `output'
statistics describing the operating characteristics of time series we care
about.''
\end{quote}
\hfill{\small(Lucas, 1980, pp.~709--710)}

\noindent Vale la pena subrayar la fecha: 1980. Mucho antes de que existiera el
software que hoy damos por sentado, el programa de investigación ya estaba
planteado en términos computacionales. Escribir el modelo como código no era un
accesorio: \emph{era} la tarea.

\subsection{La ejecución: el experimento computacional}

Quienes convirtieron esa visión en una metodología operativa fueron Finn
Kydland y Edward Prescott. Su modelo de ciclos reales ---``Time to Build and
Aggregate Fluctuations'' (1982)--- no solo propuso una teoría; propuso una
\emph{manera de trabajar} que Kydland resumiría, al recibir el Nobel, en una
frase:

\begin{quote}\itshape
``The key tool macroeconomists use is the computational experiment.''
\end{quote}
\hfill{\small(Kydland, 2004, Nobel Prize Lecture, p.~341)}

\noindent En este enfoque el modelo se trata como un instrumento de medición
que, antes de usarse, debe \emph{calibrarse}: ``In a sense, models are measuring
devices: they need to be calibrated'' (Kydland, 2004, p.~342). Calibrar
significa fijar los parámetros con información de largo plazo y de estudios
microeconómicos, de modo que el modelo reproduzca hechos que ya conocemos,
para luego interrogarlo sobre los que no.

¿Por qué calibrar, y no simplemente estimar el modelo con toda la maquinaria
econométrica? La respuesta tiene una anécdota reveladora. Thomas Sargent, que
había liderado el esfuerzo por llevar las expectativas racionales al terreno
econométrico formal, recuerda el desencanto que condujo a la calibración:

\begin{quote}\itshape
``\ldots after about five years of doing likelihood ratio tests on rational
expectations models, I recall Bob Lucas and Ed Prescott both telling me that
those tests were rejecting too many good models.''
\end{quote}
\hfill{\small(Sargent, en Evans \& Honkapohja, 2005, pp.~561--583)}

\noindent Los tests formales rechazaban modelos que, por lo demás, resultaban
útiles y económicamente sensatos. La calibración nació, en parte, como una
respuesta pragmática a esa tensión: conservar las implicaciones del modelo que
nos interesan y no exigirle que supere todas las pruebas probabilísticas a la
vez.

\begin{remark}
Esta tensión ---estimar formalmente versus calibrar--- no está zanjada, y hoy
conviven ambas tradiciones. Dynare, de hecho, hace las dos cosas: permite
calibrar y simular un modelo (lo que haremos en esta sesión) y también
estimarlo bayesianamente con datos reales (el paso siguiente, que solo
asomaremos al final). La elección es del investigador, no de la herramienta.
\end{remark}

\subsection{Las herramientas: de los algoritmos a Dynare}

La visión de Lucas y la metodología de Kydland--Prescott necesitaban, para
volverse práctica cotidiana, algoritmos concretos y software que los
implementara. Dos piezas resultaron decisivas. Por el lado de los modelos
\emph{lineales}, Blanchard y Kahn (1980) dieron la condición que determina
cuándo un modelo de expectativas racionales tiene una solución única y estable
---la misma condición que Dynare verifica hoy con el comando \code{check} y que
usaremos en esta sesión. Por el lado de los modelos \emph{no lineales} de
previsión perfecta, Fair y Taylor (1983) propusieron un método iterativo
(\emph{extended path}) que es el antecedente directo del solver determinista de
Dynare.

El software que unificó todo esto lo creó \textbf{Michel Juillard}, quien
inició el proyecto Dynare en~\textbf{1994} en el CEPREMAP (París). En su origen,
Dynare resolvía modelos deterministas de previsión perfecta mediante un método
de Newton sobre el sistema apilado; con los años incorporó la simulación
estocástica por perturbación y la estimación bayesiana, hasta convertirse en el
estándar \emph{de facto} de la macroeconomía cuantitativa. En febrero de 2025,
el Banque de France celebró sus treinta años con una conferencia dedicada.

\begin{figure}[H]
\centering
\begin{tikzpicture}[font=\footnotesize, >={Stealth[length=2mm]}]
  % línea
  \draw[thick, upblue] (0,0) -- (13.2,0);
  % hitos
  \foreach \x/\y/\a in {
      0/1976/{Crítica de Lucas},
      2.4/1980/{Lucas: ``FORTRAN program''\\ Blanchard--Kahn},
      5.1/1982/{Kydland--Prescott:\\ Time to Build},
      7.3/1983/{Fair--Taylor:\\ previsión perfecta},
      9.8/1994/{Juillard crea\\ \textbf{Dynare} (CEPREMAP)},
      12.6/2025/{30 años\\ (Banque de France)}}{
    \filldraw[upblue] (\x,0) circle (2.3pt);
    \node[upblue, font=\footnotesize\bfseries] at (\x,0.35) {\y};
    \node[align=center, text width=2.6cm, font=\scriptsize] at (\x,-0.75) {\a};
  }
\end{tikzpicture}
\caption{Evolución cuantitativa de la macroeconomía y la influencia de Dynare.}
\label{fig:timeline}
\end{figure}

\begin{remark}
El hilo conductor de esta historia es sencillo: \emph{escribir el modelo como
código no reemplaza la teoría; la completa}. La pizarra plantea el sistema de
ecuaciones; la computadora lo resuelve, lo simula y lo confronta con los datos.
Esa es, en una frase, la filosofía de diseño de Dynare: escribir el modelo casi
como se haría en el papel y dejar que la máquina haga el resto.
\end{remark}

% ===================== REFERENCIAS =====================
\begin{thebibliography}{99}
\addcontentsline{toc}{chapter}{Referencias}

\bibitem{bk1980} Blanchard, O.~J.\ \& Kahn, C.~M.\ (1980). The Solution of
Linear Difference Models under Rational Expectations. \emph{Econometrica},
48(5), 1305--1311.

\bibitem{ft1983} Fair, R.~C.\ \& Taylor, J.~B.\ (1983). Solution and Maximum
Likelihood Estimation of Dynamic Nonlinear Rational Expectations Models.
\emph{Econometrica}, 51(4), 1169--1185.

\bibitem{kp1982} Kydland, F.~E.\ \& Prescott, E.~C.\ (1982). Time to Build and
Aggregate Fluctuations. \emph{Econometrica}, 50(6), 1345--1370.

\bibitem{kp1996} Kydland, F.~E.\ \& Prescott, E.~C.\ (1996). The Computational
Experiment: An Econometric Tool. \emph{Journal of Economic Perspectives},
10(1), 69--85.

\bibitem{kydland2004} Kydland, F.~E.\ (2004). \emph{Quantitative Aggregate
Theory}. Nobel Prize Lecture, 8 de diciembre de 2004.

\bibitem{lucas1980} Lucas, R.~E.\ Jr.\ (1980). Methods and Problems in Business
Cycle Theory. \emph{Journal of Money, Credit and Banking}, 12(4), 696--715.

\bibitem{prescott1986} Prescott, E.~C.\ (1986). Theory Ahead of Business-Cycle
Measurement. \emph{Federal Reserve Bank of Minneapolis Quarterly Review},
10(4), 9--22.

\bibitem{sargent2005} Sargent, T.~J.\ (2005). An Interview with Thomas J.\
Sargent (por G.~W.\ Evans \& S.~Honkapohja). \emph{Macroeconomic Dynamics},
9, 561--583.

\bibitem{dynare} Adjemian, S., Bastani, H., Juillard, M., Karamé, F.,
Mutschler, W., Pfeifer, J., Ratto, M., Rion, N.\ \& Villemot, S.\
\emph{Dynare: Reference Manual}. CEPREMAP, Dynare Working Papers.

\end{thebibliography}

\end{document}
